\clearpage
\section{Experiments}\label{sec:experiments}

\subsection{Objectives and scope}
This chapter specifies the benchmark that was executed, using its persisted forecasts and environment records. The descriptive questions are: how does point forecast accuracy vary across tasks and model families; how are errors associated with training-series characteristics; and how sensitive is TimesFM 2.5 to context length and inference flags? These questions organize the completed analysis; they are not preregistered hypotheses. The study uses pretrained models without updating their weights, a pooled tabular formulation fitted on each task's training observations, and individually fitted statistical baselines. No cross-dataset weight adaptation, fine-tuning, probabilistic calibration or rolling-origin experiment was performed.

\subsection{Datasets and benchmark construction}\label{sec:datasets}
There are ten distinct dataset panels and eleven dataset--horizon tasks: ETTh1 is evaluated twice. Each variable is forecast as a univariate series. Melting a wide data table creates series identifiers; it does not make the forecasting problem multivariate. The tabular models pool training examples within a task, but their lag inputs still concern only the target series. The other implementations process each series separately, possibly in batches.

The evaluated versions are the cached M4 Hourly data and LongHorizon2 archive loaded with \texttt{datasetsforecast} 1.0.1. Normalization is disabled when loading registered LongHorizon2 groups. Exchange and ILI are read directly from the same archive and melted by the project loader. No dataset-wide standardization, resampling, interpolation or outlier removal is applied. Table~\ref{tab:datasets} gives the effective training lengths and forecasting horizons. Appendix~\ref{sec:inventory} gives exact timestamp boundaries and additional inventory information. These measurements were recomputed from the cached data, rather than inferred from benchmark names. The library's loading and ETT truncation are documented in its source~\cite{datasetsforecast_lh2}.

\begin{table}[!htbp]
\centering
\caption{Evaluated data (datasetsforecast 1.0.1). LH2: cached LongHorizon2 archive, unnormalized. $s/m$: configured seasonal lag / MASE lag. *Exchange elapsed time refers to the archive labels.}
\label{tab:datasets}
\small
\setlength{\tabcolsep}{3.5pt}
\begin{tabular}{lllrlll}
\hline
Task & Version & Step & Series & Train length & $H$ / duration & $s/m$ \\
\hline
M4 Hourly & M4 & h & 414 & 700--960 & 48 / 2 days & 24 / 24 \\
ETTh1/48 & LH2 & h & 7 & 14352 & 48 / 2 days & 24 / 24 \\
ETTh1/96 & LH2 & h & 7 & 14304 & 96 / 4 days & 24 / 24 \\
ETTh2/48 & LH2 & h & 7 & 14352 & 48 / 2 days & 24 / 24 \\
ETTm1/48 & LH2 & 15min & 7 & 57552 & 48 / 12 hours & 96 / 96 \\
ETTm2/48 & LH2 & 15min & 7 & 57552 & 48 / 12 hours & 96 / 96 \\
ECL & LH2 & h & 321 & 26256 & 48 / 2 days & 24 / 24 \\
Traffic & LH2 & h & 862 & 17496 & 48 / 2 days & 24 / 24 \\
Weather & LH2 & 10min & 21 & 52647 & 48 / 8 hours & 144 / 144 \\
Exchange & LH2 & D & 8 & 7540 & 48 / 48 days* & 7 / 1 \\
ILI & LH2 & W-TUE & 7 & 918 & 48 / 48 weeks & 52 / 52 \\
\hline
\end{tabular}
\end{table}


\paragraph{M4 Hourly.}
M4 is a forecasting competition archive whose public repository distributes separate training and test files~\cite{m4_data}. All 414 hourly series are retained. The loader concatenates the supplied portions, and withholding the last 48 observations recovers the supplied hourly holdout. Training lengths are 700--960, with evaluated lengths 748--1,008. The labels are integer positions beginning at one; the final training position for a series of length $N_i$ is $N_i-48$. The public numeric panel does not provide reliable calendar dates, physical units or a semantic identity for each anonymized series. Thus the domain is the heterogeneous competition collection, not an inferred application such as electricity. Padding cells in the rectangular source files are not observations; the loaded sequences have no nonfinite values or index gaps. Lag 24 represents one day, while $H=48$ represents two nominal days.

\paragraph{Electricity Transformer Temperature (ETT).}
The ETT-small release concerns two transformer stations in China. ETTh1 and ETTh2 are hourly versions; ETTm1 and ETTm2 use 15-minute observations, as confirmed by the files, despite imprecise wording about minutes in the repository description. Each has seven numeric variables: high, middle and low useful load (HUFL, MUFL, LUFL), the corresponding useless-load variables (HULL, MULL, LULL), and oil temperature (OT). The source identifies their meanings but does not specify defensible physical units for all stored columns~\cite{ett_data}. This study predicts all seven separately, including OT, without using the other six as simultaneous covariates.

The cached hourly files contain 17,420 observations per variable. The 15-minute files contain 69,680. Both begin on 2016-07-01 and extend to 2018-06-26 (19:00 and 19:45 respectively). The library retains only the first 14,400 or 57,600, ending on 2018-02-20. The holdout is taken after this truncation, so the discarded tail is not part of the test. The hourly training lengths are 14,352 for $H=48$ and 14,304 for ETTh1 at $H=96$; the quarter-hourly training length is 57,552. All four evaluated panels are regular and finite, but contain signed values. Daily seasonal lags are 24 and 96 respectively. Consequently $H=48$ covers two days in ETTh and only twelve hours in ETTm; ETTh1 at $H=96$ covers four days.

\paragraph{Electricity consumption (ECL).}
The original UCI electricity-load collection (2011--2014) contains 370 clients at 15-minute resolution during 2011--2014, with consumption represented in kW. Its documentation specifies zero-filled periods before some clients begin, and special treatment of daylight-saving transitions~\cite{electricity_data}. The long-horizon benchmark is an hourly derivative with 321 variables~\cite{lai_data}. The cached ECL version has 26,304 observations per variable and identifiers \texttt{0}--\texttt{319} plus \texttt{OT}; these are anonymous client series, not an oil-temperature target. Its labels span 2016-07-01 02:00 to 2019-07-02 01:00, inconsistent with the original source years. They must therefore be treated as archive labels. The exact hourly aggregation/unit conversion and mapping back to original clients are not recorded locally, so errors are described in stored units rather than asserted to be kWh. All values are finite and nonnegative, with regular hourly spacing. The 26,256 training observations precede a 48-hour holdout. Source zero-filling means that numerical completeness does not guarantee continuous measurement activity.

\paragraph{Traffic.}
The source describes hourly road occupancy fractions, between zero and one, from California Department of Transportation sensors on San Francisco Bay Area freeways~\cite{lai_data}. The evaluated \texttt{TrafficL} version has 862 anonymous variables, labelled \texttt{0}--\texttt{860} and \texttt{OT}, and 17,544 observations each. All are finite, nonnegative and regularly spaced. Training uses 17,496 observations, with a final 48-hour holdout and seasonal lag 24. The local labels span 2016-07-01 02:00 to 2018-07-02 01:00. The upstream description associates the collection with 2015--2016 but also contains an inconsistent month count; neither statement establishes the observation dates of this relabelled archive. Sensor locations, exclusion rules, original missing-value treatment and a sensor-ID crosswalk are unavailable in the evaluated files.

\paragraph{Weather.}
The long-horizon Weather benchmark is associated with Jena meteorological measurements during 2020 at ten-minute intervals~\cite{autoformer_paper}. The station documentation defines measured and derived variables~\cite{jena_weather}. The 21 evaluated columns comprise pressure; air, potential and dew-point temperature; relative humidity; saturation, actual and deficit vapour pressure; specific humidity; water-vapour concentration; air density; wind speed, maximum wind speed and direction; rainfall and rain duration; shortwave radiation; photosynthetically active radiation (PAR) and its maximum; logger temperature; and a final column renamed \texttt{OT}. Appendix~\ref{sec:inventory} lists the headers and units. The physical identity of that renamed final column is not independently established by the local archive.

There are 52,695 observations per variable, from 2020-01-01 00:10 to 2021-01-01 00:00. One interval jumps from 2020-05-29 09:30 to 11:10, omitting nine expected ten-minute labels. All stored values are finite; five variables contain negative values, and physical plausibility of every value has not been audited. The loader preserves the gap. Index-based models and feature extraction consequently treat the adjacent stored rows as adjacent observations, while Prophet retains their timestamps. The gap is outside the final holdout. Training has 52,647 observations; the final 48 observations run from 16:10 to midnight and span eight hours from the forecast origin. Lag 144 represents nominal daily seasonality.

\paragraph{Exchange rates.}
The source collection contains daily exchange rates for Australia, the United Kingdom, Canada, Switzerland, China, Japan, New Zealand and Singapore, described as covering 1990--2016~\cite{lai_data}. The cached derivative has eight variables (\texttt{0}--\texttt{6}, \texttt{OT}) and 7,588 rows, labelled consecutively from 1990-01-01 to 2010-10-10. Those daily labels cannot verify the source calendar or distinguish trading days from calendar days. The anonymous headers do not preserve the country mapping or quote convention; all values remain in stored exchange-rate units. No missing values or label gaps occur. Training uses 7,540 observations and testing uses 48, corresponding to 48 days on the archive clock only. The feature and statistical seasonal lag is seven, whereas the MASE scale uses lag one. This is a configured evaluation choice, not evidence that all currencies exhibit weekly seasonality.

\paragraph{Influenza-like illness (ILI).}
The long-horizon collection derives from US CDC weekly ILI surveillance, described by the benchmark authors as covering 2002--2021~\cite{autoformer_paper}. ILI surveillance measures visits meeting a symptom-based case definition, rather than laboratory-confirmed influenza alone~\cite{cdc_ili}. The exact archive contains seven related variables: weighted and unweighted ILI percentages, counts for ages 0--4 and 5--24, total ILI cases, number of providers, and \texttt{OT}. The last variable's original name and unit are not preserved; it is not safe to infer its meaning from the generic target label. The percentages and count columns have different scales and are statistically dependent. All 966 observations per variable are finite and nonnegative. The archive assigns regular Tuesdays from 2002-01-01 to 2020-06-30; these labels do not establish the original surveillance weeks or revision vintage. There are 918 training observations and a 48-week holdout, almost a year, with seasonal and MASE lag 52. This duration makes ILI qualitatively different from a 48-step sub-daily task.

\paragraph{Horizon interpretation and departures from standard protocols.}
A common 48-observation horizon simplifies the implementation, but does not equalize elapsed time, number of seasonal cycles, or difficulty. The additional ETTh1 horizon examines a longer continuation on the same dataset. Standard long-horizon evaluations commonly use chronological train/validation/test partitions and multiple forecast windows; for example, Autoformer reports 6:2:2 ETT and 7:1:2 other splits~\cite{autoformer_paper}. Here each loaded series contributes exactly one final window, almost all its preceding history is available subject to model caps, and raw-scale point metrics are reported. These results cannot be read as a reproduction of those papers' normalized multivariate MAE/MSE tables.

\subsection{Forecasting protocol and model configurations}\label{sec:protocol}
For series $i$, let $y_{i,1:N_i}$ be its sorted loaded observations and let $T_i=N_i-H$. The forecast origin is the last training observation $T_i$. Models receive observations through $T_i$ and predict $\hat y_{i,T_i+h}$ for $h=1,\ldots,H$. The test timestamps identify forecast positions; test target values are not inputs. Context means the history supplied to a model, and can be shorter than the available training series. There is no validation set, repeated seed experiment or rolling origin. The settings in Table~\ref{tab:modelsettings} describe the saved main run.

\begin{table}[!htbp]
\centering\small
\caption{Nine implemented comparators and main-run settings. History lengths are upper bounds in observations; $s$ is the configured seasonal lag.}\label{tab:modelsettings}
\begin{tabular}{>{\raggedright\arraybackslash}p{.18\linewidth}>{\raggedright\arraybackslash}p{.73\linewidth}}\hline
Model & Implementation and settings \\\hline
Linear regression & scikit-learn 1.9.0; median imputation, intercept, pooled direct horizon rows; 10,000-row cap. \\
TabPFN & \texttt{TabPFNRegressor()} from 8.0.8; same pooled rows; float32 inputs; prediction batches of 1,024; package-default checkpoint. \\
TimesFM 2.5 & \texttt{google/timesfm-2.5-200m-pytorch}; package 2.0.1; context 1,024; all five studied flags enabled; returned point output. \\
Chronos-Bolt & \texttt{amazon/chronos-bolt-base}; chronos-forecasting 1.5.3; context 2,048; batch 64; median forecast; CUDA, bfloat16 weights. \\
Prophet & 1.3.0; per-series fit; 4,320-point tail; automatic daily/weekly/yearly seasonality; four workers; uncertainty samples zero. \\
AutoARIMA & StatsForecast 2.0.3; 5,000-point tail; restricted nonseasonal approximate search (details below); four workers. \\
AutoETS & StatsForecast 2.0.3; 5,000-point tail; automatic ETS selection, season length $s$; four workers. \\
AutoTheta & StatsForecast 2.0.3; 5,000-point tail; automatic Theta selection, season length $s$; four workers. \\
SeasonalNaive & StatsForecast 2.0.3; repeat the last $s$ observations; no explicit history cap; four workers. \\\hline
\end{tabular}
\end{table}

\subsubsection{Pooled tabular formulation}
Linear regression and TabPFN use one supervised table per task. With zero-based training cutoff $c$, candidate rows pair $c=0,\ldots,T_i-2$ with $h=1,\ldots,\min(H,T_i-1-c)$. The target is $y_{i,c+h}$ in this indexing. For each series, the scale is $a_i=T_i^{-1}\sum_t|y_{i,t}|+10^{-8}$, computed over its complete training portion. Both target and lag values are divided by $a_i$; predictions are multiplied by $a_i$ afterward. This scale is train-only with respect to the final test, but includes observations later than individual internal supervised cutoffs; these rows are not separate validation origins.

The predictors are $h$, $h/H$, $\log a_i$, sine and cosine of the target observation's phase modulo $s$, and lag values $y_{i,c-\ell+1}/a_i$. Time phase uses a one-based observation index, not weekday or holiday metadata. An unavailable lag is NaN. Linear regression applies training-column median imputation before ordinary least squares; TabPFN receives the float32 table directly, including missing lag entries. There is no project-level median imputation for TabPFN. The configured lag sets are reproduced in Appendix~\ref{sec:inventory}.

Candidate counts determine proportional row allocations under a total cap of 10,000 per task. Integer allocations use floors, minimum one per nonempty series when the budget permits, and fractional-remainder redistribution. A NumPy generator with seed 42 samples without replacement within each series; sampled rows are sorted before assembly. The saved metadata reports 10,000 training rows for every task. Test rows all use the single final origin, so predictions for later horizons do not feed earlier predictions back into their lag features. This is a pooled direct multi-horizon regression, distinct from recursively fitted univariate autoregression. It also differs from TabPFN-TS v1, which constructs time-derived features and processes individual series separately~\cite{tabpfn_ts}; the comparison does not evaluate that published method itself.

\subsubsection{Pretrained forecasting models}
TimesFM receives each series' final at most 1,024 values as float32 arrays. The wrapper loads and caches one checkpoint, compiles the forecast configuration for the task, and submits the list of series to the package. Its requested maximum horizon is $H$. The returned point array is scored, with quantile outputs discarded. In the inspected installed implementation, the returned point is channel five, the median channel; it should not be assumed to be the mean merely because the API also exposes a mean channel. The package handles internal batching; the wrapper does not fix a separate TimesFM batch size.

Chronos-Bolt loads weights with the requested bfloat16 dtype on CUDA, uses float32 input contexts of at most 2,048 observations, and requests quantile 0.5 in batches of 64. Returned quantiles are converted to float32 before scoring. It sets \texttt{limit\_prediction\_length=False}; ETTh1 at 96 steps therefore uses the package's extension beyond its native 64-step output. The model card describes Bolt as a patch-based encoder--decoder predicting quantiles directly, unlike the original Chronos token-sampling model~\cite{chronos_bolt_card}. The saved per-model version string is ``unknown'', but distribution metadata records \texttt{chronos-forecasting} 1.5.3. Checkpoint identifiers are recorded without immutable revision hashes. TabPFN's package-default weight identity is less specific still. These omissions limit reconstruction of the original pretrained state.

\subsubsection{Statistical baselines and Prophet}
AutoETS and AutoTheta receive only the configured season length in addition to package defaults. AutoARIMA is deliberately constrained: seasonal terms are disabled; $p,q\leq1$, $d\leq1$, $P=Q=D=0$, maximum total order two, starting AR/MA orders zero, at most five models in the configured search, approximation enabled, and drift disabled. Although the season length is passed, seasonal ARIMA estimation is not enabled. Its results therefore characterize a computationally restricted comparator, not an exhaustive ARIMA search.

Prophet represents a series through trend and Fourier seasonal components, with optional holiday effects~\cite{prophet_paper,prophet_seasonality}. Here it is fitted separately to each series with no added holidays or regressors. Package defaults specify linear growth, additive seasonality, 25 candidate changepoints in the first 80\% of the fitting history, changepoint prior scale 0.05, seasonality prior scale 10, and MAP fitting (\texttt{mcmc\_samples=0}). Daily, weekly and yearly seasonalities are selected automatically according to history span and sampling interval; they are not all forced on. The wrapper caps the fitting history at 4,320 observations, uses four spawned workers and scores only \texttt{yhat}, with uncertainty sampling disabled. Prophet's internal scaling uses its default absolute-maximum rule on the fitting data.

For M4, integer position $t$ is mapped to 2000-01-01 plus $(t-1)$ hours solely for Prophet. This preserves the hourly ordering but does not establish real weekdays or holiday dates. A fixed shift of the time origin rotates the sine/cosine basis coefficients, so synthetic dating alone does not turn weekly seasonality into noise. For the other panels Prophet receives the archive timestamps.

Predictions are clipped below at zero when a series is classified as nonnegative from training data. The tabular, TimesFM, Chronos and StatsForecast wrappers use the complete training series for this decision. Prophet instead checks its retained fitting tail. Thus clipping can differ if a series had negative values only in older history. No test observations determine clipping. All numerical results refer to these postprocessed predictions.

\subsubsection{Reproducibility and information budgets}
Model histories are unequal by design: 1,024 and 2,048 observations for the two forecasting foundation models, 5,000 for fitted StatsForecast baselines, and 4,320 for Prophet. The tabular models can sample rows across the entire training period but retain at most 10,000 rows per task. Neither equal elapsed histories nor equal computational budgets are imposed. In particular, the tabular models share information between series through fitting, whereas the local baselines do not.

Environment records dated 2026-07-06 and 2026-07-07 include configurations, package versions, dataset counts, model status/details, Python/platform information, and an NVIDIA RTX A5000 with 24,564 MiB reported memory. They record commit \texttt{39f5848} with a dirty working tree. Current source/configuration checks support the implementation description but cannot reconstruct that uncommitted historical source state. The saved seed controls the project's row sampling; there is no verified global seeding of every model/library. Checkpoint hashes, per-fit selected statistical orders, full compiled TimesFM objects, elapsed training times and peak-memory measurements are not recorded. Persisted forecasts nevertheless permit exact re-evaluation of the reported errors.

\subsection{Metrics and aggregation}\label{sec:metrics}
For a set $A$ of forecast observations, with errors $e_{i,h}=\hat y_{i,T_i+h}-y_{i,T_i+h}$ and $n=|A|$, the implementation computes
\begin{align}
\mathrm{MAE}(A)&=\frac1n\sum_{(i,h)\in A}|e_{i,h}|, &
\mathrm{RMSE}(A)&=\sqrt{\frac1n\sum_{(i,h)\in A}e_{i,h}^{2}},\\
\mathrm{sMAPE}(A)&=\frac1n\sum_{(i,h)\in A}\frac{200|e_{i,h}|}{|y_{i,T_i+h}|+|\hat y_{i,T_i+h}|+10^{-8}}.
\end{align}
sMAPE follows the 0--200 convention (reported as percentages). Two zeros contribute zero; tiny values can make a small absolute error a large percentage error. MAE is in stored units, and RMSE emphasizes large errors. Mixing columns with different units makes their pooled raw errors difficult to interpret, even within Weather or ILI.

For training length $T_i$, use lag $m_i=m$ if $T_i>m$, and otherwise $m_i=1$. The scale and MASE are
\begin{align}
q_i&=\frac{1}{T_i-m_i}\sum_{t=m_i+1}^{T_i}|y_{i,t}-y_{i,t-m_i}|,\\
\mathrm{MASE}(A)&=\frac{1}{|A^+|}\sum_{(i,h)\in A^+}\frac{|e_{i,h}|}{q_i},\qquad
A^+=\{(i,h)\in A:q_i>0,\ q_i\text{ finite}\}.
\end{align}
If fewer than two training observations exist, the scale is NaN. Zero/nonfinite scales are excluded from MASE, not replaced with an epsilon; an empty valid set yields NaN. All saved scales in both runs are positive and finite, so no scored rows are excluded here. The lag is the configured seasonal lag except for Exchange ($m=1$). MASE provides a train-derived scale reference~\cite{mase_paper}; a value below one compares test error with an in-sample naive error, not necessarily with the actual SeasonalNaive test forecast.

The three scopes are overall (all series and steps within one task/model), series (one series across $H$), and horizon (all series at a fixed step). Overall RMSE takes the square root after pooling squared errors; it is not the mean of per-series RMSEs. With complete equal-horizon panels, overall MAE, sMAPE and MASE equal the arithmetic mean of the corresponding per-series scores. No arithmetic average of raw errors across datasets is used. Minima in result tables are selected before rounding, with exact ties sharing emphasis.

\subsection{Training-series characteristics}\label{sec:characteristics}
Features are extracted before the held-out continuation. Length counts the complete training series; the other eight features use its final at most 10,000 observations. Mean and standard deviation use population moments, and $\mathrm{CV}=\sigma/|\mu|$ is undefined when $|\mu|\leq10^{-12}$. Lag-one and seasonal autocorrelation are computed by FFT-based \texttt{statsmodels.acf}, with its default unadjusted covariance normalization~\cite{acf_docs}. For a tail $x_{1:n}$ this corresponds to
\[
 r_k=\frac{\sum_{t=k+1}^{n}(x_t-\bar x)(x_{t-k}-\bar x)}{\sum_{t=1}^{n}(x_t-\bar x)^2}.
\]
The seasonal ACF is omitted when the tail is too short ($n\leq s+1$). Constant series or $n<3$ have undefined ACF, decomposition strengths and entropy.

Robust STL with period $s$ is fitted when $s\geq2$ and $n\geq2s+1$. With decomposition $x=T+S+R$, the strengths are
\[
 F_T=\max\{0,1-\operatorname{Var}(R)/\operatorname{Var}(T+R)\},\quad
 F_S=\max\{0,1-\operatorname{Var}(R)/\operatorname{Var}(S+R)\}.
\]
A zero component-plus-remainder variance makes the associated strength undefined. These definitions follow the STL feature construction in \emph{Forecasting: Principles and Practice}~\cite{stl_features}.

For $n\geq8$, SciPy's default periodogram (constant detrending, boxcar window) estimates spectral power~\cite{periodogram_docs}. Retain its strictly positive powers $P_j$, set $p_j=P_j/\sum_jP_j$, and compute $E=-\sum_{j=1}^{K}p_j\log p_j/\log K$ when $K\geq2$. A more uniform retained spectrum gives higher entropy. A finite-sample white-noise periodogram is not perfectly uniform, so its estimated entropy need not equal one. Irregular timestamps are not explicitly modelled in these index-based features.

Per-series metrics are joined to features on task and series ID. Spearman's rank correlation is computed for each model/feature pair both within tasks and across all task-series entries, dropping missing pairs. Fewer than three pairs or a constant feature yields an undefined correlation; tied values use the rank convention of pandas/SciPy. Valid-pair counts accompany the selected results. These are descriptive associations, without confidence intervals, significance tests or multiple-comparison correction. The original analysis concerns sMAPE; an additional MASE view is derived from the same saved tables and kept separately.

\subsection{TimesFM test-set sensitivity study}\label{sec:sensitivity}
The separate run evaluates M4 Hourly, ETTh1/48, ETTm1/48 and Weather/48 using the fixed TimesFM checkpoint identifier. Six maximum contexts are requested: 512, 1,024, 2,048, 4,096, 8,192 and 16,256. Five additional variants each disable one of \texttt{normalize\_inputs}, \texttt{use\_continuous\_quantile\_head}, \texttt{force\_flip\_invariance}, \texttt{infer\_is\_positive}, and \texttt{fix\_quantile\_crossing}. The 1,024-context variant with all flags enabled is the shared baseline, making eleven variants and 44 task--variant combinations.

The installed implementation rounds the maximum horizon to a 128-point output patch and requires context plus rounded horizon not to exceed 16,384; contexts are multiples of the 32-point input patch. Thus 16,256 is the largest context in this grid for the 48-step forecast. M4 histories have at most 960 observations, and ETTh1/48 has 14,352, limiting the actual observations available at larger settings. Only ETTm1 and Weather have enough history to fill every requested context.

Metadata confirms the requested context and one changed flag in each ablation, but records the constructor arguments rather than the post-compile object. The wrapper contains a compatibility fallback that would retain only context, horizon and normalization arguments after a constructor TypeError. Inspection of the installed package supports the intended flags; historical logs do not independently prove that this fallback never occurred. Results therefore describe saved configured variants with that provenance limitation. Only point forecasts are assessed, and wrapper nonnegativity clipping remains active even when positive inference is disabled.

More broadly, inference configuration includes context/horizon budgets, normalization, sign/positivity handling, quantile processing and covariate integration via XReg. Weight adaptation is a different operation: the official repository documents a Transformers/PEFT LoRA fine-tuning example for TimesFM 2.5~\cite{timesfm_repo}. Full or partial weight updates would require an optimization and evaluation protocol beyond this study. No such updates are executed here. All variants are evaluated on the same final test set; this is test-set sensitivity analysis, not validation-selected hyperparameter optimization. The best observed configuration cannot be reported as an unbiased tuned test result.

\subsection{Comparator rationale and correspondence to prior studies}\label{sec:comparators}
Table~\ref{tab:papers} connects the implemented families to primary evaluations. It distinguishes the original papers from the checkpoint versions used here. Statistical baselines recur in the literature and provide interpretable reference forecasts. Linear regression checks the shared tabular representation, while Prophet adds a trend/seasonality model requested for examination. These choices do not amount to a full replication, and no literature scores are inserted into the experimental ranking.

\begin{table}[!htbp]\centering\footnotesize
\caption{Primary-source comparator and protocol mapping. Baseline lists are selected families, not exhaustive reimplementations.}\label{tab:papers}
\begin{tabular}{>{\raggedright\arraybackslash}p{.15\linewidth}>{\raggedright\arraybackslash}p{.23\linewidth}>{\raggedright\arraybackslash}p{.23\linewidth}>{\raggedright\arraybackslash}p{.28\linewidth}}\hline
Paper/version; model & Data/protocol and metrics & Baseline families & Correspondence and differences \\\hline
TimesFM v4 (2024); original TimesFM~\cite{timesfm_paper} & Monash, Darts, ETT; scaled MAE, ETT MAE/MSE; ETT horizons 96/192 & ETS/ARIMA, learned Monash models, LLMTime; PatchTST on ETT & Classical references overlap; this study uses 2.5, different tails/horizons and raw-scale evaluation. \\
Chronos v2 (2024); Chronos-T5\newline\cite{chronos_paper} & In-domain and 27-dataset zero-shot partition; MASE/WQL & Naive, SeasonalNaive, AutoETS/Theta/ARIMA; DeepAR, TFT, PatchTST, DLinear and others & Four statistical families shared; Bolt checkpoint and point-only scoring differ. \\
TabPFN-TS v1 (2025)~\cite{tabpfn_ts} & 24 AutoGluon-TS datasets; relative MASE geometric aggregation & SeasonalNaive, AutoETS/ARIMA/Theta, DeepAR, TFT, Chronos-Mini/Large & Statistical rationale shared; this project's pooled lag/horizon adaptation and package differ. \\
Prophet preprint (Jan. 2017)~\cite{prophet_paper} & Simulated historical forecasts; horizon-dependent percentage error & Naive, mean, ARIMA, ETS, seasonal naive, TBATS & Prophet and several families included; this study has one origin, automatic defaults and no analyst holiday adjustments. \\\hline
\end{tabular}\end{table}
